\section{Preparation for Cartan's criterion}

Consider the base field $F$. If $\operatorname{char}F=0$ then $\mathbb{Q}\subseteq F$. Let $\pi: \sum\limits_{i=1}^n \mathbb{Q}\alpha_i \to \mathbb{Q}$ be a $\mathbb{Q}$-linear map.

Consider
\[
    \psi=
    \begin{pmatrix}
        \pi(\alpha_1) & 0             & \cdots & 0             \\
        0             & \pi(\alpha_2) & \cdots & 0             \\
        \vdots        & \vdots        & \ddots & \vdots        \\
        0             & 0             & \cdots & \pi(\alpha_n)
    \end{pmatrix}.
\]
This $\psi$ is a replica of $\varphi_s$. Immediately we know that $\mathrm{ad}(\psi)$ is a replica of $\mathrm{ad}(\varphi_s)$.

\begin{lemma}
    Let $A$ and $B$ be subspaces of $\mathfrak{gl}(V)$ ($\operatorname{dim}_F(V)<\infty$). let
    \[
        M=\{ \varphi \in \mathfrak{gl}(V) \mid [\varphi,B] \subseteq A \}
    \]
    Then $M$ is a Lie subalgebra of $\mathfrak{gl}(V)$. If $\varphi \in M$, forall $m \in M$, $\operatorname{Tr}(\varphi m)=(0)$, then $\varphi$ is nilpotent.
\end{lemma}

\begin{proof}
    We shall prove that $\varphi_s=0$, hence $\varphi=\varphi_n$ is nilpotent.

    Choose a basis of $V$ such that
    \[
        \varphi_s=
        \begin{pmatrix}
            \alpha_1 & 0        & \cdots & 0        \\
            0        & \alpha_2 & \cdots & 0        \\
            \vdots   & \vdots   & \ddots & \vdots   \\
            0        & 0        & \cdots & \alpha_n
        \end{pmatrix},
        \qquad \alpha_i\in F.
    \]
    Since $\operatorname{char}F=0$, we have $\mathbb{Q}\subseteq F$. Set
    \[
        W:=\sum_{i=1}^n \mathbb{Q}\alpha_i\subseteq F.
    \]
    Let
    \[
        \pi:W\to \mathbb{Q}
    \]
    be any $\mathbb{Q}$-linear map, and define
    \[
        \psi=
        \begin{pmatrix}
            \pi(\alpha_1) & 0             & \cdots & 0             \\
            0             & \pi(\alpha_2) & \cdots & 0             \\
            \vdots        & \vdots        & \ddots & \vdots        \\
            0             & 0             & \cdots & \pi(\alpha_n)
        \end{pmatrix}.
    \]
    Then $\psi$ is a replica of $\varphi_s$.

    We claim that $\psi\in M$. Indeed, since $\varphi\in M$, we have
    \[
        [\varphi,B]\subseteq A.
    \]
    Passing to the quotient $\mathfrak{gl}(V)/A$, this means that the induced map
    \[
        \operatorname{ad}(\varphi):B+A/A\to \mathfrak{gl}(V)/A
    \]
    is zero. Hence its semisimple part
    \[
        \operatorname{ad}(\varphi_s)
    \]
    also acts as zero on $B+A/A$. Since $\operatorname{ad}(\psi)$ is a replica of $\operatorname{ad}(\varphi_s)$, it follows that $\operatorname{ad}(\psi)$ also vanishes on $B+A/A$. Therefore
    \[
        [\psi,B]\subseteq A,
    \]
    so $\psi\in M$.

    By the hypothesis of the lemma, for every $m\in M$ one has
    \[
        \operatorname{Tr}(\varphi m)=0.
    \]
    Applying this to $m=\psi$, we get
    \[
        \operatorname{Tr}(\varphi\psi)=0.
    \]
    Now
    \[
        \operatorname{Tr}(\varphi\psi)
        =
        \operatorname{Tr}(\varphi_s\psi)+\operatorname{Tr}(\varphi_n\psi).
    \]
    Since $\psi$ is a polynomial in $\varphi_s$, it commutes with $\varphi_s$, and hence also with $\varphi_n$. Therefore $\varphi_n\psi$ is nilpotent, so
    \[
        \operatorname{Tr}(\varphi_n\psi)=0.
    \]
    Consequently,
    \[
        0=\operatorname{Tr}(\varphi_s\psi)=\sum_{i=1}^n \alpha_i\,\pi(\alpha_i).
    \]
    This holds for every $\mathbb{Q}$-linear map $\pi:W\to\mathbb{Q}$.

    Now choose a $\mathbb{Q}$-basis $\beta_1,\dots,\beta_r$ of $W$, and write
    \[
        \alpha_i=\sum_{j=1}^r a_{ij}\beta_j,
        \qquad a_{ij}\in\mathbb{Q}.
    \]
    Let $\pi_k:W\to\mathbb{Q}$ be the coordinate functional defined by
    \[
        \pi_k(\beta_j)=\delta_{jk}.
    \]
    Then
    \[
        0=\sum_{i=1}^n \alpha_i\,\pi_k(\alpha_i)
        =\sum_{i=1}^n \alpha_i a_{ik}
        =\sum_{i=1}^n \left(\sum_{j=1}^r a_{ij}\beta_j\right)a_{ik}
        =\sum_{j=1}^r \left(\sum_{i=1}^n a_{ij}a_{ik}\right)\beta_j.
    \]
    Since $\beta_1,\dots,\beta_r$ are linearly independent over $\mathbb{Q}$, we obtain
    \[
        \sum_{i=1}^n a_{ij}a_{ik}=0
        \qquad
        \text{for all }1\le j,k\le r.
    \]
    In particular, taking $j=k$, we get
    \[
        \sum_{i=1}^n a_{ik}^2=0
        \qquad
        \text{for each }k.
    \]
    As each $a_{ik}\in\mathbb{Q}$, this implies
    \[
        a_{ik}=0
        \qquad
        \text{for all }i,k.
    \]
    Hence $\alpha_i=0$ for every $i$, so $\varphi_s=0$.

    Therefore
    \[
        \varphi=\varphi_n
    \]
    is nilpotent. This proves the lemma.
\end{proof}

